%==============================================================================
% CONCLUSION
%==============================================================================
\chapter*{Conclusion}
\addcontentsline{toc}{chapter}{Conclusion}
\markboth{Conclusion}{Conclusion}

\section*{Summary of the work}

This research project has resulted in a working system, \textbf{AutoOFFBEAT},
which makes it possible to drive the fuel performance code OFFBEAT from requests
expressed in natural language. The agent creates computational cases from
validated templates, runs the solver, attempts to repair certain crashes
automatically, post-processes the results and evaluates safety margins against
design criteria.

The system was then extended towards a fuel rod digital twin: monitoring of
safety criteria, prognosis of threshold crossing, dashboard, and a
Gaussian-process surrogate returning the results of a simulation within a few
milliseconds, with an associated uncertainty and a fidelity verified by
cross-validation.

A reference simulation over a two-year irradiation cycle made it possible to
verify that the complete chain reproduces the expected physical behaviour, in
particular the closure of the pellet-cladding gap and the associated switch of
the cladding stress from compression to tension. The surrogate, rebuilt over a
realistic irradiation domain, predicts the date of that closure to within
\SI{4.4}{\percent} for a power it has never encountered, in \SI{2}{\milli\second}
instead of \SI{36}{\second}.

The second strand of the work, less expected at the outset, is the
\textbf{measured stress-testing of the system}. A fault-injection benchmark
quantified the actual reach of self-healing; an evaluation of the decision layer
measured, then improved by eleven points, the accuracy of tool selection. Both
benchmarks are re-runnable and delivered with the code: they constitute a
non-regression harness for the rest of the project.

\section*{What I learned}

Three lessons seem to me worth retaining.

The first is a matter of physics. I was discovering fuel performance entirely at
the start of the internship. Understanding that the pellet-cladding gap governs
both the thermal and the mechanical behaviour of the rod, and being able to read
the signature of contact in simulation curves, is the clearest technical gain.

The second concerns the software architecture of systems built on language
models. The principle that structures the whole project — the model orchestrates,
it does not compute; deterministic first, model as fallback — carries well beyond
OFFBEAT.

The third is methodological, and is perhaps the most important. The most
dangerous defects encountered were not crashes but wrong results that did not
look wrong. No automatic tool flagged them; they were found by confronting a
numerical value with a physical expectation. A safety computation chain is not
validated by checking that it runs, but by checking that what it produces makes
sense.

This third point admits a corollary I had not anticipated: a validation metric
can itself be misleading. The surrogate displayed a coefficient of determination
above $0.998$, obtained by a rigorous method, over a domain that did not cover
the real problem. Checking that a model is well fitted and checking that it is
fitted to the right question are two distinct operations, and only the second
falls to physical judgement.

\section*{Outlook}

Several concrete follow-ups stand out, in order of priority.

\textbf{Validate the safety thresholds.} The criteria currently written into the
knowledge base are marked as unvalidated. Confirming them with the supervisor
and the literature is the prerequisite to any serious use of the safety analysis.

\textbf{Settle the scope.} Moving from the single rod to the assembly changes the
nature of the problem and remains an open question.

\textbf{Shift the effort from repair to validation.} The analysis of chapter~4
shows that validating inputs before execution is more reliable than enriching the
base of crash patterns. A first geometric validation was added; it could be
extended systematically to all the exposed parameters.

\textbf{Extend the surrogate.} The surrogate covers only two parameters. Since
its reconstruction over a realistic domain showed that a design of twenty-five
simulations fits in half an hour, adding a third variable — the initial gap
geometry, for instance — is now within reach. The code is written for an
arbitrary number of variables.

\textbf{Make time-step self-healing reliable.} The benchmark showed that the fix
is undone by the solver's adaptive controller; acting on the maximum time step
rather than on the initial one is an identified and localised correction.

\textbf{Complete the validation of the thresholds.} The correlation for the
decrease of the melting temperature with burnup was settled at the end of the
internship by returning to the literature: the value carried by the project came
from the MATPRO correlation, since superseded in the reference codes
(appendix~\ref{ann:safetykb}). The same approach remains to be carried out on the
three other criteria, including the yield strength of Zircaloy, strongly
temperature-dependent and for the moment given as a single value.

\textbf{Extend the tool-selection benchmark.} The current twenty-six requests
were enough to reveal a major ambiguity between two descriptions. A broader set,
and above all multi-step requests, would allow the chaining of tools to be
evaluated rather than the first call alone.

\clearpage
